\section{Notation and function-space setting}
\label{app:notation}
% =====================================================================

Throughout, $\Omega = [-1,1]^2$ is identified with the flat torus
$\mathbb{T}^2 = (\R/2\Z)^2$, so that the spectral operators below are exact
rather than approximate. Let $L^2(\Omega;\R^C)$ carry the inner product
$\langle u,v\rangle = \int_\Omega u(x)^\top v(x)\,\dd x$ and norm
$\|u\| = \langle u,u\rangle^{1/2}$. The Fourier transform is
\begin{equation}
\hat{u}(\vk) = \tfrac{1}{|\Omega|}\int_\Omega u(x)\,e^{-\mathrm{i}\vk\cdot x}\dd x,
\qquad
u(x) = \sum_{\vk\in\Lambda} \hat{u}(\vk)\,e^{\mathrm{i}\vk\cdot x},
\qquad
\Lambda = (\pi\Z)^2 ,
\label{eq:fourier}
\end{equation}
with Parseval identity $\|u\|^2 = |\Omega|\sum_{\vk\in\Lambda}|\hat u(\vk)|^2$.
The lowest non-zero wavenumber on $\Omega$ is $|\vk|_{\min} = \pi$.

Let $\mathcal{S}_{++}^2$ denote the cone of symmetric positive-definite
$2\times2$ matrices. A \emph{diffusion tensor field} is a tuple
$\mD = (\mD_1,\dots,\mD_C) \in (\mathcal{S}_{++}^2)^C$, acting channelwise. The
\emph{reaction field} is a map $f_\theta \in C^1(\R^C;\R^C)$ applied pointwise,
$(f_\theta u)(x) = f_\theta(u(x))$.

\begin{definition}[Morphogen operator]
\label{def:operator}
For $\mD \in (\mathcal{S}_{++}^2)^C$ and $f_\theta \in C^1(\R^C;\R^C)$ the
\emph{morphogen operator} is the semilinear evolution equation
\begin{equation}
\partial_t \vz = \mathcal{A}_{\mD}\,\vz + f_\theta(\vz),
\qquad
(\mathcal{A}_{\mD}\vz)_c := \nabla\!\cdot\!\big(\mD_c \nabla \vz_c\big),
\label{eq:operator}
\end{equation}
on $\Omega$ with periodic boundary conditions and initial datum
$\vz(\cdot,0)=\vz_0 \in L^2$.
\end{definition}

Equation~\eqref{eq:operator} is the classical semilinear parabolic form: a
strongly elliptic second-order generator plus a globally Lipschitz, bounded
perturbation. Existence and uniqueness of a mild solution on $[0,\infty)$ is
standard, following from Propositions~\ref{prop:contraction} and~\ref{prop:reaction}.

% =====================================================================
\subsection*{Spectral representation}
\label{app:spectral}

\begin{lemma}[Fourier symbol]
\label{lem:symbol}
For each channel $c$, $\mathcal{A}_{\mD_c}$ is diagonal in the basis
\eqref{eq:fourier} with symbol
\begin{equation}
\widehat{\mathcal{A}_{\mD_c} u}(\vk) \;=\; -\,q_c(\vk)\,\hat{u}(\vk),
\qquad
q_c(\vk) := \vk^\top \mD_c \vk \;\ge\; \lambda_{\min}(\mD_c)\,|\vk|^2 \;\ge\; 0 .
\label{eq:symbol}
\end{equation}
\end{lemma}

\begin{proof}
Writing $\mD_c=(d_{ij})$ and $u = e^{\mathrm{i}\vk\cdot x}$,
$\nabla u = \mathrm{i}\vk\,u$, hence $\mD_c\nabla u = \mathrm{i}\,\mD_c\vk\,u$
and $\nabla\!\cdot\!(\mD_c \nabla u) = \mathrm{i}\vk \cdot (\mathrm{i}\mD_c\vk)u
= -(\vk^\top \mD_c\vk)\,u$. The inequality is the Rayleigh bound for symmetric
$\mD_c$.
\end{proof}

Consequently the diffusion semigroup admits the closed form used in the
implementation,
\begin{equation}
\big(e^{\tau \mathcal{A}_{\mD}}u\big)^{\wedge}(\vk) \;=\;
  e^{-q_c(\vk)\tau}\,\hat{u}_c(\vk),
\label{eq:semigroup}
\end{equation}
which is evaluated by a pair of FFTs at cost $O(CHW\log HW)$ per half-step. No
spatial discretization of $\nabla\!\cdot\!(\mD\nabla\cdot)$ is ever formed, so the
diffusive update carries no truncation error beyond the spectral representation
of $\vz$ itself.

\begin{remark}
Isotropic $\mD_c=d_cI$ gives a rotation-invariant symbol $d_c|\vk|^2$; the
general case has elliptical level sets aligned with the eigenvectors of $\mD_c$,
permitting different decay lengths along and across an anatomical axis. The
isotropy ablation of Section~\ref{sec:ablations} measures whether this freedom is
used.
\end{remark}

% =====================================================================
